\documentclass{article}
\usepackage[T1]{fontenc}
\usepackage{lmodern}
\usepackage{graphicx}
\usepackage{amsmath,amssymb}
\usepackage[a4paper,margin=2cm]{geometry}
\usepackage[absolute,overlay]{textpos}
\setlength{\TPHorizModule}{1mm}
\setlength{\TPVertModule}{1mm}
\usepackage[indonesian]{babel}
\usepackage{hyperref}
\setlength{\emergencystretch}{2em}
% unit: Halaman 27

\begin{document}

% === Halaman 27 (python/native) ===

27

B. Vigènere Cipher 

• Vigenere Cipher termasuk ke dalam kelompok cipher abjad-majemuk 

(polyalpabetic cipher ). Ini berbeda dengan cipher abjad-tunggal 

(monolphabetic cipher).

• Perbedaan polyalphabetic cipher dengan monoalphabetic cipher:

➢  Cipher abjad-tunggal: satu kunci untuk semua huruf plainteks 

       Contoh: Caesar cipher

➢

Cipher abjad-majemuk: setiap huruf menggunakan  kunci berbeda. 

       Contoh: Vigenere cipher

• Cipher abjad-majemuk dibuat untuk mengatasi kelemahan cipher abjad-

Tunggal yang mudah dipecahkan dengan teknik analisis frekuensi

\end{document}
